\section{Discussion}\label{sec:discussion}

\subsection{Comparison with Literature: Ramanujan Adjacent Graphs}\label{sec:lit_comparison}

A systematic search---arXiv, Google Scholar, Wikipedia---found \textbf{zero} prior results on Ramanujan properties of adjoint commutation graphs of Lie algebras. The Sunada equivalence between spectral Ramanujan and zeta Ramanujan is well-known for \emph{regular} graphs, but the adjoint graphs are inherently \emph{irregular} (Cartan hub + root vertices). Our complete failure of Sunada's equivalence for all adjoint graphs (except su(2)) is a genuine new result.

The connection with West, Bringewatt, and Brady's quantum chaos framework~\cite{west2025quantum} is particularly significant: Ramanujan algebras are maximally chaotic (negative average OTOC, $2\times$ faster mixing). The Standard Model sits deep in the optimally chaotic regime.

\subsection{Comparison with Connes: Spectral Action vs. Evolution Graph}\label{sec:connes_comparison}

Connes' spectral action program~\cite{connes1994noncommutative,chamseddine1997spectral} chooses the finite geometry $F$ to \emph{reproduce} the Standard Model. SS-PEG takes the Cartan-Weyl graph as the \emph{input}, not the output. The distinction is fundamental:

\begin{itemize}
    \item \textbf{Connes:} Imposed symmetry $\to$ spectral action $\to$ SM
    \item \textbf{SS-PEG:} Relational substrate $\to$ Killing minimization $\to$ emergent algebra $\to$ graph $\to$ SM
\end{itemize}

SS-PEG demonstrates \emph{emergence of symmetry} rather than \emph{imposition of symmetry}. The Cartan-Weyl graph is not chosen to match the SM---it is computed from the variational principle.

\subsection{Theoretical Implications}\label{sec:implications}

\begin{itemize}
    \item \textbf{Gauge symmetry is emergent, not fundamental.} The graph topology emerges from the Killing spectrum, which emerges from the variational principle, which emerges from the relational substrate.
    
    \item \textbf{Spacetime is algebraic, not arena.} The spacetime hierarchy (Levels 0--4) is a sequence of graph topological transitions, not a progression of background geometries.
    
    \item \textbf{The universe's topology is a layered composition of all five archetypes.} Internal gauge cycles (SU(3), SU(2)) + abelian strand (U(1)) $\to$ spacetime backbone (so(3,1)) $\to$ conformal structure (su(2,2)).
\end{itemize}

\subsection{Limitations and Open Questions}\label{sec:limitations}

\begin{itemize}
    \item \textbf{Continuum limit ($d \to \infty$):} Open. The Ramanujan divergence for classical families suggests that large-rank algebras cannot crystallize from the variational principle.
    
    \item \textbf{Connecting Paper 2 with Paper 3:} Deriving SM parameters from graph invariants. The five topological building blocks provide the bridge.
    
    \item \textbf{Tensegrity equation:} The functional form $g(\Kil_\la)$ remains to be determined.
\end{itemize}

\subsection{Connection to Paper 3: The Top of the Chain}\label{sec:paper3_connection}

Paper 3 takes the five spectral blocks from Paper 2:
\begin{itemize}
    \item $\mathcal{R}_2$, $\mathcal{R}_3$, $R_{EE}$: spectral ratios
    \item $h_2^{\vee 4} = 16$, $h_3^{\vee 4} = 81$: Coxeter amplifiers
\end{itemize}

and uses them to predict $\alpha$, $\alpha_s$, $\alpha_2$, fermion masses, CKM, and PMNS parameters with $0$ free parameters. The $73$ predictions tested yield $72/73$ within $1\%$ error.

The graph does not merely classify---it \emph{computes} observables.
